\documentclass[11pt]{scrartcl}

\usepackage[top=1.5cm]{geometry}
\usepackage{url}
\usepackage{float}
\usepackage{listings}
\usepackage{xcolor}
\usepackage{graphicx}
\usepackage{booktabs}

\setlength{\parindent}{0em}
\setlength{\parskip}{0.5em}

\renewcommand{\thesubsection}{\arabic{subsection}}

\lstdefinestyle{dmrsql}{
  language=SQL,
  basicstyle=\small\ttfamily,
  keywordstyle=\color{magenta!75!black},
  stringstyle=\color{green!50!black},
  showspaces=false,
  showstringspaces=false,
  commentstyle=\color{gray}}

\lstdefinestyle{dmrPython}{
  language=Python,
  basicstyle=\small\ttfamily,
  keywordstyle=\color{magenta!75!black},
  stringstyle=\color{green!50!black},
  showspaces=false,
  showstringspaces=false,
  commentstyle=\color{gray}}

\title{
  \textbf{\large Projektaufgabe 3 } \\
  Phase 1 – Das EDGE Modell (4 P) \\
  {\large Datenmanagement jenseits von Relationen}
}

\author{
  Gruppen Nummer (A17) \\
  \large Acikyol Aleyna, 12032843 \\
  \large Özcan Ahmet, 12311287 \\
  \large Sadykova Nargiz, 12129337
}

\begin{document}

\maketitle\thispagestyle{empty}

Dieses Reporting Template dient der Vorbereitung der Abgabe von Phase 1.

\subsection*{EDGE Modell Import des Toy Beispiels (1 Punkte)}

Beschreiben Sie kurz, welches Verfahren (inkl. Programmiersprache) Sie für den Import verwenden. Gab es Probleme beim Implementieren des Parsers? (0.5 P).

Für den XML-Import wurde Python mit der Standardbibliothek \texttt{xml.etree.ElementTree} verwendet.
Das Dokument wird als String eingelesen, Sonderzeichen (Umlaute, \texttt{\&szlig;} etc.) werden vor dem
Parsen per \texttt{str.replace()} normalisiert, da die DTD-Entities von \texttt{ElementTree} nicht
automatisch aufgelöst werden. Anschließend wird der Baum in eine eigene \texttt{Node}-Klasse überführt
und per Tiefensuche (\emph{DFS}) in die Datenbank geschrieben.

Probleme gab es ausschließlich mit den HTML-Entities in den DBLP-Daten (z.\,B. \texttt{\&uuml;}),
die \texttt{ElementTree} ohne DTD nicht kennt. Die Vorverarbeitung per \texttt{normalize\_input()}
löst dieses Problem zuverlässig.

Zeigen Sie (z.B. mittels Screenshot des Konsolenoutputs), dass Ihr Parser alle Elemente (und deren Inhalt) des Toy-Beispiels korrekt erkennt. (0.5 P).

Der folgende Ausschnitt zeigt den Konsolenoutput von \texttt{phase1\_import.py} für das Toy-Beispiel
(4 Publikationen):

\begin{lstlisting}[style=dmrPython, frame=single]
=== EDGE-Modell Toy-Beispiel (Konsolenausgabe) ===
[bib] bib
  [venue] vldb
    [year] vldb_2023
      [article] SchmittKAMM23
        [author]: Daniel Ulrich Schmitt
        [author]: Daniel Kocher
        [author]: Nikolaus Augsten
        [author]: Willi Mann
        [author]: Alexander Miller
        [title]: A Two-Level Signature Scheme for Stable Set Similarity Joins.
        [pages]: 2686-2698
        [year]: 2023
        [volume]: 16
        [journal]: Proc. VLDB Endow.
        [number]: 11
        [ee]: https://www.vldb.org/pvldb/vol16/p2686-schmitt.pdf
        [ee]: https://doi.org/10.14778/3611479.3611480
        [url]: db/journals/pvldb/pvldb16.html#SchmittKAMM23
      [article] SchalerHS23
        [author]: Christine Schaeler
        [author]: Thomas Huetter
        [author]: Martin Schaeler
        [title]: Benchmarking the Utility of w-event Differential Privacy ...
        [pages]: 1830-1842
        [year]: 2023
        [volume]: 16
        [journal]: Proc. VLDB Endow.
        [number]: 8
        [ee]: https://www.vldb.org/pvldb/vol16/p1830-schaler.pdf
        [ee]: https://doi.org/10.14778/3594512.3594515
        [url]: db/journals/pvldb/pvldb16.html#SchalerHS23
  [venue] sigmod
    [year] sigmod_2022
      [inproceedings] HutterAK0L22
        [author]: Thomas Huetter
        [author]: Nikolaus Augsten
        [author]: Christoph M. Kirsch
        [author]: Michael J. Carey 0001
        [author]: Chen Li 0001
        [title]: JEDI: These aren't the JSON documents you're looking for?
        [pages]: 1584-1597
        [year]: 2022
        [booktitle]: SIGMOD Conference
        [ee]: https://doi.org/10.1145/3514221.3517850
        [crossref]: conf/sigmod/2022
        [url]: db/conf/sigmod/sigmod2022.html#HutterAK0L22
  [venue] pacmmod
    [year] pacmmod_2023
      [article] ThielKAHMS23
        [author]: Konstantin Emil Thiel
        [author]: Daniel Kocher
        [author]: Nikolaus Augsten
        [author]: Thomas Huetter
        [author]: Willi Mann
        [author]: Daniel Ulrich Schmitt
        [title]: FINEX: A Fast Index for Exact & Flexible Density-Based ...
        [pages]: 71:1-71:25
        [year]: 2023
        [volume]: 1
        [journal]: Proc. ACM Manag. Data
        [number]: 1
        [ee]: https://doi.org/10.1145/3588925
        [url]: db/journals/pacmmod/pacmmod1.html#ThielKAHMS23
\end{lstlisting}

Alle Elemente und deren Inhalte werden korrekt erkannt. Die Attribute \texttt{mdate} und \texttt{orcid}
werden wie gefordert ignoriert.

Geben Sie die Definition Ihrer Tabellen im Edge Modell an. Zeigen Sie (einen Ausschnitt) des Toy Beispiels nach dem Import als Edge Model in die DB (1 P).

Das EDGE-Modell besteht aus zwei Tabellen. Die \textbf{Node}-Tabelle speichert alle Knoten mit einem
numerischen Primärschlüssel (\texttt{id}), einem symbolischen Bezeichner (\texttt{s\_id}, z.\,B.
Publikationsschlüssel oder Venue-Name), dem Knotentyp (\texttt{type}) sowie dem Textinhalt
(\texttt{content}, nur für Blattknoten). Die \textbf{Edge}-Tabelle hält alle gerichteten Eltern-Kind-Kanten
als Paare (\texttt{from}, \texttt{to}).

\begin{lstlisting}[style=dmrsql, frame=single]
CREATE TABLE node (
    id      SERIAL PRIMARY KEY,
    s_id    VARCHAR,
    type    VARCHAR NOT NULL,
    content TEXT
);

CREATE TABLE edge (
    "from"  INTEGER NOT NULL REFERENCES node(id),
    "to"    INTEGER NOT NULL REFERENCES node(id),
    PRIMARY KEY ("from", "to")
);
\end{lstlisting}

Ausschnitt des importierten Toy-Beispiels (erste Publikation, analog zu Tabelle 1 der Aufgabenstellung):

\begin{table}[H]
\centering
\begin{tabular}{r l l l | r r}
\toprule
\multicolumn{4}{c|}{\textbf{Node}} & \multicolumn{2}{c}{\textbf{Edge}} \\
\texttt{id} & \texttt{s\_id} & \texttt{type} & \texttt{content} & \texttt{from} & \texttt{to} \\
\midrule
0 & bib            & bib          & null                  & 0 & 1 \\
1 & vldb           & venue        & null                  & 1 & 2 \\
2 & vldb\_2023     & year         & null                  & 2 & 3 \\
3 & SchmittKAMM23  & article      & null                  & 3 & 4 \\
4 & null           & author       & Daniel Ulrich Schmitt & 3 & 5 \\
5 & null           & author       & Daniel Kocher         & 3 & 6 \\
6 & null           & author       & Nikolaus Augsten      & 3 & 7 \\
7 & null           & author       & Willi Mann            & 3 & 8 \\
8 & null           & author       & Alexander Miller      & 3 & 9 \\
9 & null           & title        & A Two-Level Sig. [...] & 3 & 10 \\
\ldots & \ldots & \ldots & \ldots & \ldots & \ldots \\
\bottomrule
\end{tabular}
\caption{Ausschnitt des Toy-Beispiels nach Import ins EDGE-Modell}
\end{table}

\subsection*{Berechnung der XPath Achsen im EDGE Modell (2 Punkte)}

Die vier XPath-Achsen (\texttt{ancestor}, \texttt{descendant}, \texttt{following-sibling},
\texttt{preceding-sibling}) werden jeweils als rekursive SQL-Anfrage (\texttt{WITH RECURSIVE})
in der Datenbank berechnet. Python dient ausschließlich als Controller, der die Anfrage absetzt
und das Ergebnis ausgibt.

Geben Sie für jede der geforderten Achsen, die Knotenmenge des Ergebnisses, sowie die Anzahl der Ergebnisknoten aus.

\begin{table}[h]
  \centering
    \begin{center}
      \begin{tabular}{ l | c c }
        \toprule
        Achse & Ergebnisknoten IDs & Ergebnisgröße\\
        \midrule
        ancestor (Daniel Ulrich Schmitt, id=5)  & 4, 3, 2, 1          & 4  \\
        descendants (vldb\_2023, id=3)          & 4--31               & 28 \\
        following SchmittKAMM23 (id=4)          & 19 (SchalerHS23)    & 1  \\
        preceeding SchmittKAMM23 (id=4)         & (leer)              & 0  \\
        following SchalerHS23 (id=19)           & (leer)              & 0  \\
        preceeding SchalerHS23 (id=19)          & 4 (SchmittKAMM23)   & 1  \\
        \bottomrule
      \end{tabular}
    \end{center}
  \caption{Ergebnisse der XPath-Berechnung}
  \label{tab:ErgebnisseDerXPathBerechnug}
\end{table}

Sofern Sie Annahmen zum Aufbau der Daten getroffen haben, geben Sie diese nachfolgend an:
\begin{itemize}
  \item Die Dokumentreihenfolge wird durch die aufsteigenden \texttt{node.id}-Werte abgebildet,
        da alle Knoten in DFS-Reihenfolge (Vorbestellung) in die Datenbank eingefügt werden.
        \texttt{following-sibling} entspricht damit Geschwisterknoten mit \texttt{id > v\_id},
        \texttt{preceding-sibling} entsprechend \texttt{id < v\_id}.
  \item Jeder Knoten hat höchstens einen Elternknoten (Baumstruktur), sodass die Elternsuche
        in der \texttt{edge}-Tabelle eindeutig ist.
  \item Die Attribute \texttt{mdate} und \texttt{orcid} werden beim Import ignoriert.
\end{itemize}

\subsection*{Zeitmamagement}

Benötigte Zeit pro Person (nur Phase 1): \textbf{ca. 3 Stunden}

\end{document}
